\section{Conclusion}\label{concl}
In this paper, we presented an exhaustive token-based evaluation of MAL in the verbal domain on the UD collection of 180 languages. We proposed a classification based on a metric which we defined in order to estimate a MAL effect for any given language satisfying a given threshold of occurrences for relevent configurations. We obtained a reliable classification of 131 languages, well-distributed between IE and non-IE languages, as well as between VO and OV languages. This classification allowed us to confirm MAL as a typologically wide-spread  preference across the languages of the world, but importantly, it showed that MAL is not an absolute universal principle. Not only we found a number of languages in which MAL is absent or trivial, but we also found languages on the strict oppoite side of the cline, thus showing a clear MAL dispreference.

The second contribution of our study is the finding that contrary to what has been assumed previously, MAL does not apply equally on the right and the left side of the verb. MAL is stronger on the right side of the verb than on the left side of the verb, while Anti MAL is stronger on the left side. This is a crutial finding since it implies that length-based effects are not necessarily similar in the pre- and postverbal domains, against the widespeard assumption of the contrary. Not only our results show a clear difference between the pre- and postverbal domains, but they also show that VO and OV languages have different preference profiles in these domains. VO languages are more likely to show a strong MAL preference in the postverbal doamin, while  OV languages are more likely to show a strong MAL preference in the preverbal domain. Importantly, similar observations on the differences between pre- and postverbal domains have previsously made been on length-based ordering preferences in inindivual languages such as Persian \citep{faghiriGlossa2020}. Moreover, there are alternative psycholinguistic production-based accounts of mirror-image word order preferences (SbL and LbS) in VO and OV languages, which posit differences between the two domains \citep[for an overview, see][]{faghiriGlossa2020}. The latter namely claim that verbs exert strong influence over the constituents that follow them and, accordingly, assume that formal constraints are more strict in the preverbal domain than in the postverbal domain \citep{StallingsEtal}. Our study brings a fresh light into this ongoing debate.
 